%--
\chapter{计算模型}

%----------------------------------------
\section{计算模型的引入}

我们首先不严格地讨论算法的直观概念。

算法是``the step-by-step procedures that computers use to solve problems"
    \footnote{这一说法源自于Quanta Magazine的文章《Computer scientists establish the best way to traverse a graph》，
    作者Ben Brubaker，链接：
    https://www.quantamagazine.org/computer-scientists-establish-the-best-way-to-traverse-a-graph-20241025/。}
    。

我们可以借用“类比”（analog）来解释算法相关的概念
    \footnote{这里参考了Erik Demaine教授在课堂上的讲授。课程信息：
    MIT open courseware, Introduction to algorithms (6.006, fall 2011)。
    链接：https://courses.csail.mit.edu/6.006/fall11/notes.shtml。
    最后访问日期：2026.8.5。}
    。
假设我们了解一台计算机及其上运行的程序，
%
那么我们的“算法”，就是“计算机程序在数学世界中的类比”（mathematical analog of a computer program）。
%
我们描述算法所用的“伪码”，就是“计算机编程语言在数学世界中类比”（mathematical analog of a programming language）。
这里所说的伪码是指面向算法描述的，结构化的自然语言，可以是中文、英文等。
%
我们的“计算模型”，就是“现实世界中的计算机在数学世界中的类比”（mathematical analog of a computer）。

%--
\begin{table}[htbp]
    \centering
    \caption{概念的对比}
    \label{t:analog}
    %--
    \begin{tabular}{c|c}
        \toprule
        \textbf{数学世界的概念} & \textbf{计算机世界的概念} \\
        \midrule
        算法 & 程序 \\
        伪码语言 & 编程语言 \\
        计算模型 & 计算机 \\
        \bottomrule
    \end{tabular}
\end{table}


%
%%--
%\begin{table}[ht]
%    \centering
%    %
%    \begin{tabular}{|>{\centering\arraybackslash}p{5cm}|>{\centering\arraybackslash}p{5cm}|}  % 两列都居中，有竖线
%        \hline                  % 水平线
%        数学世界的概念 & 计算机世界的概念 \\        
%        \hline \hline                  % 水平线
%        算法 & 计算机程序 \\ 
%        \hline        
%        伪码 & 编程语言 \\ 
%        \hline        
%        计算模型 & 计算机 \\ 
%        \hline
%    \end{tabular}
%    %
%\end{table}

%----------------------------------------
\section{抽象的算法设计}

%----------------------------------------
\section{抽象的算法分析}